\makeatletter
\def\input@path{{./}{../}{../../}{../../../}{../../../../}{preamble/}{../preamble/}{../../preamble/}{../../../preamble/}}
\makeatother
\input{preamble_exams.tex}
\pagestyle{empty}
\newcommand{\answerlines}{\par\vspace{3mm}\hrule\vspace{3mm}\hrule\vspace{2mm}}
\newcommand{\choices}{\begin{center}\begin{tabular}{lcc}\toprule Criterion & Alternative(s) & Deciding quantity and value\\\midrule Maximax & \rule{25mm}{0.4pt} & \rule{45mm}{0.4pt}\\ Maximin & \rule{25mm}{0.4pt} & \rule{45mm}{0.4pt}\\ Minimax Regret & \rule{25mm}{0.4pt} & \rule{45mm}{0.4pt}\\ EMV/Bayes & \rule{25mm}{0.4pt} & \rule{45mm}{0.4pt}\\\bottomrule\end{tabular}\end{center}}
\begin{document}
\newgeometry{left=0.4in,right=0.4in,top=0.35in,bottom=0.35in}
\begingroup\renewcommand{\baselinestretch}{1}\small
\setlength{\parskip}{3pt}\renewcommand{\arraystretch}{1.05}
\begin{center}{\large\bfseries Grade 10 -- Term 1: C8 Classroom Activity}\end{center}
Name: \rule{75mm}{0.4pt}\hfill Date: \rule{30mm}{0.4pt}

\textbf{Goal:} Compare decision criteria, their information needs, priorities, and limitations.
All entries are profits; larger is better. Use supplied summaries unless asked to update them. Report all ties. Give brief explanations; use your notebook for extra working.

\textbf{Guided Problems}

\textbf{1. Delivery expansion.} Profits are in thousand US dollars.
\begin{center}\begin{tabular}{lrrr|rrrr}\toprule
&\multicolumn{3}{c|}{Market payoffs}&\multicolumn{4}{c}{Supplied summary}\\
Plan & Strong & Stable & Weak & Row max. & Row min. & Max. regret & EMV\\
Probability & .20 & .60 & .20 &&&&\\\midrule
$A_1$: Large fleet &100&20&-20&100&-20&55&28\\
$A_2$: Flexible fleet &75&50&10&75&10&25&47\\
$A_3$: Current fleet &40&40&35&40&35&60&39\\
$A_4$: Partners &60&65&20&65&20&40&55\\\bottomrule
\end{tabular}\end{center}
(a) Complete the selections. Maximax: highest row maximum; Maximin: highest row minimum; Minimax Regret: smallest maximum regret; EMV: highest expected profit.
\choices
(b) Complete: only \rule{25mm}{0.4pt} needs probabilities. Regret compares plans within each \rule{25mm}{0.4pt}. Match these priorities to criteria: greatest upside; best worst outcome; least maximum missed opportunity; highest expected money.
\answerlines
(c) Compare $A_2$ and $A_4$: give a numerical reason for each choice and one limitation of each criterion used. Does EMV 55 guarantee profit 55?
\answerlines

\textbf{2. Festival catering.} Profits are in thousand US dollars.
\begin{center}\begin{tabular}{lrr|rrrr}\toprule
& High & Low & Row max. & Row min. & Max. regret & EMV\\
Probability & .30 & .70 &&&&\\\midrule
$A_1$: Large contract &80&10&80&10&25&31\\
$A_2$: Small contract &45&35&45&35&35&38\\\bottomrule
\end{tabular}\end{center}
(a) Record all four recommendations and deciding values, in the same order as Problem 1.
\answerlines
(b) Low attendance: $A_1$ earns 10 and forgoes $35-10=\underline{\hspace{12mm}}$. Explain why this regret is not a cash loss. Does the regret choice protect the highest minimum profit?
\answerlines
(c) New probabilities are $(.60,.40)$. Complete $EMV(A_1)=.60(80)+.40(10)=\underline{\hspace{12mm}}$ and $EMV(A_2)=.60(45)+.40(35)=\underline{\hspace{12mm}}$. State the new EMV choice and explain why the other three choices stay unchanged.
\answerlines
\newpage
\textbf{Standard Problems}

For each problem, report all four recommendations with their deciding quantities and values. Justify comparisons numerically; mark EMV unavailable if necessary.

\textbf{3. Retail plans: agreement and ties.} Profits are in thousand CNY.
\begin{center}\begin{tabular}{lrrr|rrrr}\toprule
& Strong & Normal & Weak & Row max. & Row min. & Max. regret & EMV\\
Probability & .25 & .50 & .25 &&&&\\\midrule
$A_1$ &60&45&30&60&30&0&45\\
$A_2$ &50&40&20&50&20&10&37.5\\
$A_3$ &60&35&30&60&30&10&40\\\bottomrule
\end{tabular}\end{center}
(a) Report recommendations, including ties. (b) Compare $A_1$ with each other plan state by state. Explain why two criteria still permit ties and why agreement does not make the criteria interchangeable.
\choices\answerlines

\textbf{4. Urban farm: subsidy and constraint.} Profits are in thousand US dollars.
\begin{center}\begin{tabular}{lrrr|rrrr}\toprule
& Strong & Stable & Weak & Row max. & Row min. & Max. regret & EMV\\
Probability & .50 & .30 & .20 &&&&\\\midrule
$A_1$: Specialty &80&40&-10&80&-10&45&50\\
$A_2$: Balanced &60&55&20&60&20&20&50.5\\
$A_3$: Staple &40&45&35&45&35&40&40.5\\\bottomrule
\end{tabular}\end{center}
(a) Report unrestricted recommendations. (b) Profit must be at least 20 in every state. A subsidy adds an integer $s$, $0\leq s\leq20$, only to $A_2$'s weak-state payoff. Find the smallest $s$ making Maximin strictly prefer $A_2$ to $A_3$. Identify eligible plans. At that subsidy, calculate updated maxima, minima, regrets, and EMVs using eligible plans only, then report all four restricted recommendations.
\answerlines\answerlines

\textbf{5. Coastal ferry: no forecast.} Profits are in thousand US dollars. No reliable state probabilities are available.
\begin{center}\begin{tabular}{lrrr|rrr}\toprule
& High & Normal & Low & Row max. & Row min. & Max. regret\\\midrule
$A_1$: Two ferries &100&35&-25&100&-25&55\\
$A_2$: One ferry &75&55&10&75&10&25\\
$A_3$: Lease capacity &55&48&30&55&30&45\\
$A_4$: Current fleet &30&30&30&30&30&70\\\bottomrule
\end{tabular}\end{center}
(a) Report recommendations and explain why equal probabilities cannot be assumed. (b) Choose a criterion for an operator prioritizing missed opportunities. Contrast it with worst-case protection, giving numerical evidence and a limitation of each. (c) A later forecast gives $(.20,.30,.50)$. Calculate all EMVs and select a plan. Explain EMV's risk orientation and one limitation.
\answerlines\answerlines
\endgroup
% Restoring geometry starts the solutions on a new page.
\restoregeometry
\renewcommand{\baselinestretch}{1}\normalsize
\section*{Solutions}
\subsection*{1. Delivery expansion}
\textbf{(a)} Compare each supplied summary column:
\begin{enumerate}[label=\roman*.]
\item Maximax: $\max(100,75,40,65)=100$: $A_1$, row maximum 100.
\item Maximin: $\max(-20,10,35,20)=35$: $A_3$, row minimum 35.
\item Minimax Regret: $\min(55,25,60,40)=25$: $A_2$, maximum regret 25.
\item EMV/Bayes: $\max(28,47,39,55)=55$: $A_4$, EMV 55. No ties occur.
\end{enumerate}
\textbf{(b)} The blanks are EMV/Bayes and state. Greatest upside: optimistic Maximax; best worst outcome: conservative Maximin; least maximum missed opportunity: Minimax Regret; highest expected money: risk-neutral EMV/Bayes. Maximax and Maximin need row payoffs; regret needs the full payoff table for comparisons within states. These three require no probabilities. EMV also requires nonnegative state probabilities summing to one, here $.20+.60+.20=1$.

\textbf{(c)} Minimax Regret favors $A_2$: maximum regret $25<40$ for $A_4$. It limits missed opportunities but does not protect the highest minimum profit: $A_2$ can earn only 10. EMV favors $A_4$: $.20(60)+.60(65)+.20(20)=12+39+4=55>47$ for $A_2$. It depends on the forecast and does not guarantee 55: actual profits are 60, 65, or 20. Maximax's optimistic choice ignores $A_1$'s possible loss of 20; Maximin protects 35 with $A_3$ but ignores the upside of 100 available from $A_1$.

\subsection*{2. Festival catering}
\textbf{(a)} Maximax: $A_1$, row maximum $80>45$. Maximin: $A_2$, row minimum $35>10$. Minimax Regret: $A_1$, maximum regret $25<35$. EMV: $A_2$, EMV $38>31$. There are no ties.

\textbf{(b)} Low-state best profit is 35, so regret is $35-10=25$. The large contract earns positive profit 10; regret 25 is forgone profit, not a cash loss. State bests $(80,35)$ give regret rows $(0,25)$ and $(35,0)$. Minimax Regret chooses $A_1$ even though its minimum profit is 10, below $A_2$'s 35; it does not maximize the minimum profit.

\textbf{(c)} $EMV(A_1)=48+4=52$; $EMV(A_2)=27+14=41$. The probabilities sum to one. EMV now selects $A_1$ (52), instead of $A_2$ (38 under the original forecast). Payoffs and regrets are unchanged: Maximax stays $A_1$ (80), Maximin stays $A_2$ (35), and Minimax Regret stays $A_1$ (25).

\subsection*{3. Retail plans}
\textbf{(a)} Maximax compares $(60,50,60)$: $A_1,A_3$ tie at row maximum 60. Maximin compares $(30,20,30)$: $A_1,A_3$ tie at row minimum 30. Minimax Regret compares $(0,10,10)$: $A_1$ wins with maximum regret 0. EMV compares $(45,37.5,40)$: $A_1$ wins with EMV 45.

\textbf{(b)} $A_1$ beats $A_2$ in every state: $60>50$, $45>40$, $30>20$. Against $A_3$, it is equal in strong and weak states and better in normal: $60=60$, $45>35$, $30=30$. Maximax and Maximin ignore the normal-state difference because the row maxima and minima tie. State bests $(60,45,30)$ give regret rows $(0,0,0)$, $(10,5,10)$, and $(0,10,0)$, distinguishing $A_1$ from $A_3$. EMVs are
\begin{align*}
EMV(A_1)&=.25(60)+.50(45)+.25(30)=15+22.5+7.5=45,\\
EMV(A_2)&=.25(50)+.50(40)+.25(20)=12.5+20+5=37.5,\\
EMV(A_3)&=.25(60)+.50(35)+.25(30)=15+17.5+7.5=40.
\end{align*}
Agreement does not make the criteria interchangeable. They prioritize optimistic upside, conservative minimum, limited missed opportunities, and risk-neutral expected money, respectively. Only EMV needs probabilities. Maximax and Maximin miss the normal-state advantage here; regret ignores likelihoods; EMV depends on probability estimates and does not guarantee its average (45 versus a possible 30 for $A_1$).

\subsection*{4. Urban farm}
\textbf{(a)} Comparing supplied columns gives Maximax: $A_1$ (row maximum 80); Maximin: $A_3$ (row minimum 35); Minimax Regret: $A_2$ (maximum regret 20); EMV: $A_2$ (50.5). There are no ties.

\textbf{(b)} $A_1$ is ineligible since $-10<20$. For $0\leq s\leq20$, $A_2$ has minimum $20+s\leq40<55<60$, and $A_3$ has minimum 35; both are eligible. Strict Maximin preference requires
\[
20+s>35\iff s>15.
\]
The smallest integer is $s=16$, or \$16,000. At $s=15$ there is only a tie. At $s=16$, eligible payoff rows are $(60,55,36)$ and $(40,45,35)$. Their maxima are 60 and 45, and minima 36 and 35. State bests are $(60,55,36)$. Subtract each eligible payoff from its state best:
\[
r(A_2)=(0,0,0),\qquad r(A_3)=(20,10,1).
\]
Maximum regrets are 0 and 20. The updated EMVs are
\begin{align*}
EMV(A_2)&=.50(60)+.30(55)+.20(36)=30+16.5+7.2=53.7,\\
EMV(A_3)&=.50(40)+.30(45)+.20(35)=20+13.5+7=40.5.
\end{align*}
All four restricted criteria select $A_2$: Maximax with row maximum 60; Maximin with row minimum 36; Minimax Regret with maximum regret 0; EMV with 53.7. Maximax changes because $A_1$ is excluded; Maximin changes because the subsidy raises $A_2$'s minimum above 35. Regrets must be recomputed after excluding $A_1$.

\subsection*{5. Coastal ferry}
\textbf{(a)} Maximax: $A_1$, row maximum 100. Maximin: $A_3,A_4$, tied at row minimum 30. Minimax Regret: $A_2$, maximum regret 25. EMV and its overall recommendation are unavailable without state probabilities. Although $A_4$ has constant payoff 30, the other plans' expected values cannot be compared without a distribution. Equal probabilities add an unsupported assumption of equally likely states.

\textbf{(b)} Minimax Regret fits the missed-opportunity priority: $A_2$ limits maximum regret to 25, versus 45 for $A_3$ and 70 for $A_4$. It compares payoffs within states without probabilities, but can recommend a minimum profit of only 10. Maximin protects a minimum of 30 with $A_3$ or $A_4$, also without probabilities. It ignores upside and likelihoods: $A_3$ can earn 55 while $A_4$ never exceeds 30, yet they tie.

\textbf{(c)} All probabilities are nonnegative and $.20+.30+.50=1$.
\begin{align*}
EMV(A_1)&=.20(100)+.30(35)+.50(-25)=20+10.5-12.5=18,\\
EMV(A_2)&=.20(75)+.30(55)+.50(10)=15+16.5+5=36.5,\\
EMV(A_3)&=.20(55)+.30(48)+.50(30)=11+14.4+15=40.4,\\
EMV(A_4)&=.20(30)+.30(30)+.50(30)=6+9+15=30.
\end{align*}
EMV selects $A_3$ (40.4), with no tie. It is risk-neutral about money and depends on reliable probabilities; 40.4 is not guaranteed since $A_3$ may earn 30. The three probability-free recommendations stay unchanged because payoffs are unchanged. EMV and Maximin agree on $A_3$ for different reasons.
\end{document}


